\documentclass[11pt]{article}
\usepackage[margin=1in]{geometry}
\usepackage{amsmath,amssymb,amsthm}
\usepackage{hyperref}
\hypersetup{hidelinks}

\newtheorem{theorem}{Theorem}
\newtheorem{lemma}[theorem]{Lemma}
\newtheorem{corollary}[theorem]{Corollary}
\newtheorem{proposition}{Proposition}
\theoremstyle{definition}
\newtheorem{definition}{Definition}
\newtheorem*{remark}{Remark}

\setlength{\parskip}{4pt plus 1pt minus 1pt}
\setlength{\parindent}{0pt}

\title{Graph Structure Is Necessary for\\Information Preservation Under Bounded Context}
\author{Patrick D.\ McCarthy}
\date{}

\begin{document}
\maketitle

\begin{abstract}
Any system that (i) accumulates propositions about a world it cannot
fully observe, (ii) operates under a bounded active context, a
hard limit on the number of propositions simultaneously available for
inference, (iii) must preserve information as the proposition
count grows beyond that bound, and (iv) receives evidence that can
be corrupted or invalidated, necessarily maintains a directed graph
over its propositions, with typed edges emerging under competitive
pressure (selective necessity). We prove this claim from two
independent arguments. The \emph{information-preservation argument}
shows that the only space-creating operation that does not lose
information is factoring, extracting shared structure into named
nodes with typed edges, and that factoring \emph{is} graph
construction. The \emph{retention-dynamics argument} shows that any
representation supporting sub-linear dependency queries, contextual
retrieval, and selective removal with cascade necessarily encodes
directed adjacency, which is a semantic graph regardless of substrate.
We show that continuous-weight mechanisms (e.g., transformer
self-attention) implement graph structure rather than constitute an
alternative to it. The result applies to any system that stores
propositions under bounded context, including artificial neural
networks, institutional knowledge systems, and, to the extent that
biological representations admit propositional decomposition, biological
neural networks.
\end{abstract}

% ------------------------------------------------------------------
\section{Introduction}
\label{sec:intro}
% ------------------------------------------------------------------

Across 73{,}593 cancer patients, the number of distinct
organizational systems disrupted by a tumor's mutations predicts
survival better than the total number of mutations
\cite{McCarthy2026e}. The organizational function a gene serves, not
its molecular identity, determines prognosis. This requires that
biological information be structured by function, not stored as a
flat list. The present paper proves that this structure is necessary.

As a knowledge base grows, the capacity to reason over it, at any
single moment, does not grow proportionally. Biological working memory holds
roughly $4 \pm 1$ items \cite{Cowan2001} (Miller's classic
$7 \pm 2$ \cite{Miller1956} reflects rehearsal-augmented span,
not raw capacity). Transformer context windows,
though large, degrade in inference quality as they fill with content
irrelevant to the current problem \cite{Liu2024}. Institutional
decision makers cannot hold an entire organization's knowledge in a
single meeting. In each case, there is a hard or effective bound on
\emph{active context}: the number of propositions simultaneously
available for inference.

This paper asks a narrow question: \textbf{what are the requirements of a knowledge representation if information
is not to be lost as the number of stored propositions exceeds the
active context bound?}

The answer is a directed graph with typed edges. We prove this from two
independent arguments, each sufficient on its own:

\begin{enumerate}
\item \textbf{Information preservation} (Proposition~\ref{thm:graph-info}).
   When a bounded system must make space for new propositions, the only
   operation that creates space without losing information is
   \emph{factoring}: extracting shared structure into a named node with
   typed edges. Factoring is graph construction. Discarding without
   adjacency information is blind; lossy compression without adjacency is
   blind discarding with extra steps.

\item \textbf{Retention dynamics} (Theorem~\ref{thm:graph-dynamics}).
   Any representation that supports sub-linear dependency queries,
   answering ``what depends on this proposition?'' in less than $O(n)$
   time, necessarily stores directed adjacency. Selective removal
   with cascade (invalidating a proposition and propagating to its
   dependents) requires typed edges to scope correctly. These are the
   operations any system performs when maintaining knowledge under
   continuous evidence arrival.
\end{enumerate}

The result is substrate independent. It does not prescribe an
implementation (adjacency lists, edge tables, attention masks, synaptic
connections). It proves that the \emph{information content} of any
adequate representation is a directed graph, regardless of how that
graph is physically realized.

\subsection{Related work}

Graph-structured knowledge representations are widely used in practice:
knowledge graphs \cite{Hogan2021}, Bayesian networks \cite{Pearl1988},
description logics \cite{BaaderCalvaneseEtAl2003}, and the
RDF/OWL stack \cite{BernersLee2001} all represent knowledge as directed
labeled graphs. Guarino's formal ontology \cite{Guarino1998} grounds
this choice in the structure of the domain: entities, relations, and
constraints map naturally to nodes, edges, and types. These demonstrate that graph structure is \emph{effective} for knowledge
representation. The question of whether it is \emph{necessary}, as
opposed to convenient, has received less formal attention.

The present paper provides the missing argument: under bounded active
context, graph structure is not a design preference but a
mathematical consequence of information preservation
(Proposition~\ref{thm:graph-info}) and sub-linear maintenance
(Theorem~\ref{thm:graph-dynamics}). The typed directed edges of
description logics and OWL are not an engineering convenience but the
minimal structure any system must maintain. This result retroactively
grounds decades of practice in a formal necessity argument.

The information bottleneck principle \cite{Tishby2000} establishes that
optimal representations compress input while preserving information
relevant to output. Our information preservation argument
(Proposition~\ref{thm:graph-info}) addresses a different question: not what
to compress, but what structure the representation must have to compress
without information loss.

Simon's bounded rationality \cite{Simon1955} establishes that cognitive
agents satisfice rather than optimize. Our bounded context constraint
($C_n$) formalizes a specific consequence: when the bound is hard,
structure is not optional but necessary.

Codd's relational model \cite{Codd1970} established that normalizing
a database, extracting repeated data into separate tables linked by
foreign keys, eliminates redundancy while preserving information.
The information-preservation argument (Proposition~\ref{thm:graph-info})
derives the same structural operation from first principles: factoring
shared propositions into named nodes with typed edges is normalization.
The contribution is not the operation itself but the proof that it is
the \emph{unique} space-creating operation that preserves information
under bounded context, and that it necessarily produces graph structure.

% ------------------------------------------------------------------
\section{Formal Framework}
\label{sec:framework}
% ------------------------------------------------------------------

We introduce the minimal definitions required for the two main results.

\begin{definition}[World State Space]
\label{def:world}
Let $W$ be a measurable space of possible world states. At time $t$, the
true state $w_t \in W$ is not directly observable by any entity.
\end{definition}

\begin{definition}[Proposition and Knowledge Base]
\label{def:knowledge}
A \emph{proposition} $p$ is a claim about $W$ with an associated
likelihood function $\ell_p : W \to [0,1]$. A \emph{knowledge base} $K$
is a collection of propositions. The posterior over world states given
$K$ is:
\[
P(w \mid K) \propto P_0(w) \cdot L_K(w)
\]
where $P_0$ is a prior and $L_K$ is the joint likelihood determined by
$K$. Propositions in $K$ may have arbitrary dependencies, shared
structure, evidential interactions, causal relationships, so $L_K$ need
not factor as a product of individual likelihoods. (When propositions
are conditionally independent given $w$, $L_K(w) = \prod_{p \in K}
\ell_p(w)$; this is a special case, not assumed.) The existence of
dependencies between propositions is precisely what motivates graph
structure: if propositions were truly independent, dependency tracking
would be unnecessary.

The \emph{uncertainty} of $K$ about $w$ is the
surprisal $U(w, K) = -\log P(w \mid K)$ \cite{Shannon1948, Cover2006}.
\end{definition}

\begin{definition}[Bounded Active Context]
\label{def:context}
For an entity $n$, let $C_n$ denote its \emph{active context capacity}:
the maximum number of propositions simultaneously available for
inference. Let
$\mathrm{ctx}(t) \subseteq K$ be the set of propositions loaded into
active context at time $t$, with $|\mathrm{ctx}(t)| \leq C_n$.

The \emph{inference quality} for problem $f$ is the relevance density:
\[
\rho(f) = \frac{|K^{[f]}|}{|\mathrm{ctx}(t)|}
\]
where $K^{[f]} \subseteq K$ is the set of propositions relevant to $f$.
$\rho = 1$ when active context holds exactly the relevant propositions;
$\rho \to 0$ when context is dominated by irrelevant content.
\end{definition}

\begin{definition}[Growth and Curation Regimes]
\label{def:regimes}
$K$ has a maximum capacity $K_{max}$ determined by substrate constraints.
When $|K| < K_{max}$, $K$ is in the \emph{growth regime}: propositions
accumulate freely. When $|K| = K_{max}$, $K$ is in the \emph{curation
regime}: adding a new proposition requires making space.
\end{definition}

\begin{definition}[Contextual Retrieval]
\label{def:retrieval}
Given problem $f$, \emph{contextual retrieval} is the operation that
loads $K^{[f]}$ into active context. The retrieval cost is the number of
propositions examined to identify $K^{[f]}$.
\end{definition}

\begin{remark}[Necessity via Selection]
\label{rem:necessity}
The results in this paper show that graph structure is
\emph{cost-optimal} for information preservation under bounded context.
A reviewer may ask: is it \emph{necessary}, or merely sufficient? The
answer depends on the competitive regime. In isolation, an entity could
use any representation and survive, at higher cost. Under competition
for finite resources, the cost differential becomes a survival
differential: an entity paying $\Omega(n)$ per dependency query loses
resources to a competitor paying $O(d)$ for the same query. Under
sustained competition, cost-optimal structure is selected for, and
alternatives are eliminated. Throughout this series, ``necessary'' means
necessary for persistence under competition---selective necessity, not
logical impossibility of alternatives.
\end{remark}

% ------------------------------------------------------------------
\section{Graph Necessity from Information Preservation}
\label{sec:info}
% ------------------------------------------------------------------

\begin{definition}[Factoring]
\label{def:factoring}
Let $p_1, \ldots, p_m \in K$ share mutual information: there exists a
random variable $F^*$ (a function of $W$) such that
$I(F^*; \ell_{p_i}) > 0$ for all $i$, and the joint likelihood
decomposes as
\[
\ell_{p_i}(w) = g_i\!\bigl(\ell_{p_i}^{\text{res}}(w),\; F^*(w)\bigr)
\]
for some residual $\ell_{p_i}^{\text{res}}$ and reconstruction function
$g_i$. The decomposition need not be unique; any valid decomposition
suffices. In the general case, $g_i$ may be an arbitrary measurable
function: the definition does not require multiplicative,
additive, or any other restricted functional form. The permissiveness
is deliberate: the theorem's force comes not from restricting $g_i$
but from the requirement that $g_i$ be \emph{stored} (as a typed edge)
to enable reconstruction. The weaker the constraint on $g_i$, the
stronger the conclusion that typed edges are necessary.
\emph{Factoring} is the operation that extracts $F^*$ into a single
named, reusable unit and replaces the shared component in each $p_i$
with a link to $F^*$.

This definition is deliberately minimal. It specifies only that a
shared component is extracted and linked. It does not specify that
the links are directed, that they carry type information, or that
the result has graph structure. Proposition~\ref{thm:graph-info}
derives these properties from the requirement of posterior
preservation: direction and typing are consequences of
reconstruction, not stipulations of the definition.

\textbf{Remark.} Factoring is not a novel operation introduced by this
paper. It is the central operation of description logics
\cite{BaaderCalvaneseEtAl2003}: the TBox/ABox separation extracts shared
conceptual structure (the TBox) from individual assertions (the ABox),
and concept subsumption hierarchies are literally shared components
referenced by multiple individuals. Upper ontologies such as BFO
\cite{Arp2015} and DOLCE \cite{Masolo2003} carry this further:
approximately 35 categories in BFO organize all of biomedical knowledge,
demonstrating that real-world knowledge bases have abundant shared
structure amenable to factoring. The definition above formalizes what
these traditions have practiced for decades.
\end{definition}

\begin{definition}[Information Preservation]
\label{def:preservation}
A \emph{space-creating operation} $f : K \to K'$ reduces active context
load: $\mathrm{load}(K') < \mathrm{load}(K)$, where load counts the
propositions simultaneously required in active context. (Node count may
increase, since factoring adds a shared node, but context load decreases
because redundant copies are replaced by shared references.) Such an
operation is \emph{information-preserving} if the posterior over world
states is unchanged:
\[
P(w \mid K') = P(w \mid K) \quad \text{for all } w \in W.
\]
This is posterior preservation: the operation may reorganize, compress,
or restructure propositions, but the joint constraint that $K$ places
on $W$ is not weakened.
\end{definition}

\begin{proposition}[Graph Necessity from Information Preservation]
\label{thm:graph-info}
Let $K$ be a knowledge base with bounded active context $C_n$
(Definition~\ref{def:context}) in the curation regime
(Definition~\ref{def:regimes}). Any information-preserving
space-creating operation on $K$ (Definition~\ref{def:preservation}) as $|K|$
grows beyond $C_n$ necessarily produces named units with directed,
typed references between them---i.e., a directed graph with typed edges.
\end{proposition}

\textbf{Argument.}
In the curation regime, adding any new proposition $p_{new}$ requires
making space. Any space-creating operation $f : K \to K'$ with
$\mathrm{load}(K') < \mathrm{load}(K)$ falls into exactly one of two categories: either it
preserves the posterior (Definition~\ref{def:preservation}), or it does
not. We show that every information-preserving option produces graph
structure, and that non-preserving options require graph structure to
avoid uncontrolled information loss.

\textbf{Option A: Discard an existing proposition.}
Without explicit relational structure, the entity has no mechanism to
determine which existing propositions are redundant. Two propositions
$p_1, p_2$ that share a latent factor appear as independent entries;
there is no stored evidence that they share structure. The choice of
which to discard is therefore made without information about relational
dependencies. Discarding $p_{old}$ loses the information it encodes, and
potentially loses the information that $p_{old}$ has any relation to
other propositions, a second-order loss. This is information loss.

A natural objection: per-proposition metadata (access frequency, recency,
embedding similarity to recent queries) can inform the discard decision
without stored adjacency. This is true for a single proposition in
isolation. But propositions share structure, and metadata must reflect
that sharing to be accurate: the recency of $p_1$ is relevant to the
discard decision for $p_2$ when both ground the same action. Sharing
metadata across propositions requires storing which propositions share
which metadata, which is storing directed relationships between entries.
The metadata-based alternative converges to adjacency the moment it
must account for shared structure, which is the non-degenerate case
(Remark on shared structure, below).

\textbf{Option B: Factor: extract shared structure
(Definition~\ref{def:factoring}).}
If propositions share a common informational component $F^*$, factoring
extracts $F^*$ as a single named node with typed edges, replacing
redundant copies with shared references. Space is created without loss:
the original information is reconstructable from $F^*$ plus the
residuals. Factoring also resolves the update problem: when new evidence
revises $F^*$, the single node is updated once, and the revision
propagates to all propositions that reference it via edges:
$O(1)$ instead of $O(|K|)$.

By Definition~\ref{def:factoring}, factoring produces nodes and directed
typed edges. This is graph construction.

\textbf{Option C: Reorganization: rewrite propositions into a more
compact equivalent encoding.}
Propositions $p_1, \ldots, p_m$ could be replaced by $q_1, \ldots, q_k$
($k < m$) that are jointly equivalent: the same posterior over $W$ is
induced. This is a change of basis in the information-theoretic sense.

If the reorganization preserves all dependency relationships,
which propositions ground which, then the result is an isomorphic
knowledge base with fewer nodes. But ``preserving dependency
relationships'' requires storing those relationships, which is storing
adjacency. If the reorganization does \emph{not} preserve dependency
relationships, then downstream propositions that depended on the
original $p_i$ are silently broken: they reference nodes that no longer
exist, with no forwarding information. The entity cannot determine what
broke without adjacency.

There is a special case: reorganization that extracts shared structure
across the $p_i$ into shared nodes referenced by the $q_j$. This is
factoring (Option~B). Reorganization that does not extract shared
structure but merely rewrites individual propositions more compactly
does not create space from shared structure; it creates space from
per-proposition compression, which is lossless compression (Option~D).

A stronger form of reorganization is \emph{ontology change}: replacing
the proposition vocabulary entirely. An entity with 100 propositions
about individual animals might reconceptualize to 10 propositions about
species. The new propositions are not factored components of the old
ones; they are a different representational vocabulary. However,
ontology change does not escape the theorem. It is a one-time
translation that produces a new knowledge base $K'$. The new $K'$
still describes a structured world, still operates under bounded
$C_n$, and still accumulates propositions as evidence arrives. As
$|K'|$ grows, the propositions in $K'$ develop shared structure (the
species propositions share habitat, physiology, evolutionary
lineage), and the same curation pressure applies. Ontology change is
analogous to random reseeding in optimization: it selects a different
starting point but does not change the landscape. The landscape is
bounded context over shared-structure propositions, and every
ontology that describes a structured world converges to the same
factoring pressure.

\textbf{Option D: Lossless compression: encode propositions more
compactly without information loss.}
Lossless compression reduces the physical footprint of stored
propositions, and any system should compress to the extent possible.
But compression has a floor: information has an irreducible size
(Shannon entropy, Kolmogorov complexity). Once that floor is reached,
the active context bound $C_n$ remains. New propositions still arrive.
The bound is a constraint on how many propositions can be
\emph{simultaneously active for inference}, not on how compactly they
are encoded in storage. A compressed bitstring must be decompressed
before any proposition can be accessed or any dependency traced; the
decompressed form faces the same curation pressure as the original.
Compression is necessary but does not solve the structural problem
the proposition addresses.

\textbf{Option E: Lossy compression: replace propositions with a
summary that occupies less space.}
Compressing $p_1$ and $p_2$ into a summary $p_s$ that discards detail is
a form of discarding: the lost detail may ground other propositions or
actions. If $p_1$ has dependents, propositions that reference it,
then replacing $p_1$ with $p_s$ silently invalidates those dependents
unless the entity knows what depends on the lost detail. But knowing
what depends on $p_1$ \emph{is} knowing $p_1$'s adjacency: its
edges. Without stored adjacency, lossy compression is blind discarding:
the entity cannot evaluate what breaks. With stored adjacency, the
entity can scope the damage, notify dependents, and compress safely;
but storing adjacency is maintaining a graph. Lossy compression that
preserves correctness requires the same relational structure as
factoring. Lossy compression that ignores relational structure is
Option~A with extra steps.

\textbf{Exhaustiveness.} The partition above is complete because any
$f : K \to K'$ with $|K'| < |K|$ either preserves the posterior
(Definition~\ref{def:preservation}) or does not. If it does not, it is
lossy (Options A, E). If it does, it must produce $K'$ with fewer
propositions but the same joint likelihood. The ways to achieve this
are: extract shared structure into reusable units (B), rewrite to an
equivalent encoding with fewer propositions (C), or compress the
physical representation without changing propositional count (D). These
three exhaust the posterior-preserving options because any such
operation either reduces the number of distinct likelihood factors
(B, C) or reduces per-factor storage cost (D); there is no fourth
category, since the only degrees of freedom are how many factors
there are and how each is stored. Of these, D is orthogonal (it
addresses storage density, not the active-context bound). C either
reduces to B (when shared structure is extracted) or requires stored
adjacency to preserve dependencies. Only B (factoring) creates space,
preserves the posterior, and solves the structural problem.

\textbf{On the limits of exhaustiveness.}
The partition above is argued, not proved in the formal sense, because
the space of possible operations on a knowledge base is not itself
formally bounded. We cannot enumerate all conceivable data structures
or operations. However, this limitation does not weaken the overall
result, because Theorem~\ref{thm:graph-dynamics} provides an
independent argument that does not rely on exhaustiveness.

The falsifiable content of the theorem is not ``the representation is a
graph,'' which risks being unfalsifiable if \emph{graph} is defined
broadly enough to include any structure that works. The falsifiable
content is the \emph{cost prediction}: Theorem~\ref{thm:graph-dynamics}
establishes that answering ``what depends on $p$?''\ costs $\Omega(n)$
without stored directed adjacency and $O(d)$ with it. This is a
measurable, empirical claim. A system that achieves sub-linear
dependency query without storing any form of directed relationship
between propositions---not adjacency lists, not sparse matrices, not
learned attention patterns, not any structure from which pairwise
dependency can be read in $o(n)$ time---would falsify
Theorem~\ref{thm:graph-dynamics}. The theorem does not say ``everything
is a graph.'' It says there is a specific cost for not storing directed
adjacency, and that cost is $\Omega(n)$ per query.

\textbf{Posterior preservation forces directed, typed structure.}
Definition~\ref{def:factoring} specifies only that $F^*$ is extracted
and linked. We now derive the structural properties that posterior
preservation (Definition~\ref{def:preservation}) requires of these
links.

\emph{Nodes.} The named units---shared components $F^*$ and residual
propositions $p_i^{\text{res}}$---are the nodes. Each is an
identifiable, individually addressable unit of information.

\emph{Direction is required by reconstruction.}
To recover the original likelihood $\ell_{p_i}(w) =
g_i(\ell_{p_i}^{\text{res}}(w), F^*(w))$, the system must know that
$p_i$ depends on $F^*$---not vice versa. This asymmetry is not
stipulated but entailed: revising $F^*$ invalidates every $p_i$'s
reconstruction and requires re-evaluation; revising $p_i$ does not
invalidate $F^*$. A system that stores only symmetric (undirected)
links cannot determine which units require re-evaluation when one
unit changes---it must re-evaluate in both directions, degrading
to worst-case cascade. The reconstruction path \emph{is} the
direction: from shared component to dependent proposition.

\emph{Typing is required by correct reconstruction.}
Knowing that $p_i$ depends on $F^*$ is insufficient for
reconstruction. The system must also store \emph{how} $F^*$ enters
the reconstruction: the function $g_i$. Different propositions
depend on the same $F^*$ in different ways---$F^*$ may \emph{ground}
$p_1$ (truth value depends on $F^*$), \emph{evidence} $p_2$
(likelihood is informed by $F^*$), or \emph{cause} $p_3$ (state is
produced by $F^*$). Without the stored $g_i$, the system cannot
recover $\ell_{p_i}$ from $F^*$ and the residual, and posterior
preservation fails. The reconstruction function $g_i$ is the edge
type: it distinguishes the \emph{kind} of dependency. This is not a
design choice but a requirement---without it, the factoring is
information-losing.

\emph{Typing also scopes cascade.} When $F^*$ is invalidated,
typed edges enable cascade to follow only dependency-bearing
links, keeping cost proportional to $|D_{F^*}|$ (the true dependency
set) rather than $|\mathrm{adj}(F^*)|$ (all neighbors). This is a
consequence of having the type, not an additional requirement.

A collection of named units with directed, typed links between them
is a directed graph with typed edges, regardless of whether it is
implemented as adjacency lists, pointers, sparse matrices, or any other
medium. The graph structure is derived from the reconstruction
requirement, not stipulated in the definition of factoring.

\textbf{Remark on shared structure.} The proposition requires that
propositions share informational components. This is not an additional
assumption but a consequence of the knowledge base describing a
structured world. The proposition ``fire is hot'' is relevant to
cooking, camping, metallurgy, and fire safety; any knowledge base
spanning these domains contains this shared component. A knowledge base
with \emph{zero} shared structure would be a collection of completely
independent facts about completely unrelated domains: each proposition
grounded in its own private ontology with no overlap. Such a knowledge
base faces no curation pressure (nothing is redundant, so nothing can
be factored) and requires no structural optimization. The proposition
addresses the non-degenerate case: knowledge bases about structured
worlds, which is every knowledge base that admits inference across
domains. The success of upper ontologies confirms that this is the
typical case: BFO \cite{Arp2015} organizes all of biomedical science
under approximately 35 shared categories; DOLCE \cite{Masolo2003}
provides cross-domain foundational categories reused across engineering,
geography, and law. These are not design choices but reflections of
shared structure that already exists in the domains they describe.

\textbf{Remark on compression.} Lossless compression changes how
compactly propositions are stored, not how many can be simultaneously
active. The bounded-context premise ($|K| > C_n$) is a constraint on
active inference capacity, not storage density. No encoding scheme removes this bound. Compression is necessary
but insufficient: it delays the curation problem without eliminating
it.

% ------------------------------------------------------------------
\section{Graph Necessity from Retention Dynamics}
\label{sec:dynamics}
% ------------------------------------------------------------------

We first derive the operations that any knowledge maintenance system
must support, then prove that any representation supporting them
necessarily encodes directed adjacency.

\begin{lemma}[Necessary Maintenance Operations]
\label{lem:ops}
Let $K$ be a knowledge base under bounded active context $C_n$
(Definition~\ref{def:context}), with continuous evidence arrival, where
evidence can be corrupted or invalidated. Then $K$ must support the
following four operations:
\begin{enumerate}
\item[(i)] \textbf{Addition with provenance.}
\item[(ii)] \textbf{Contextual retrieval.}
\item[(iii)] \textbf{Dependency query.}
\item[(iv)] \textbf{Selective removal with cascade.}
\end{enumerate}
\end{lemma}

\begin{proof}
None of the three premises (bounded context, evidence arrival,
corruptibility) mentions adjacency or graph structure. We derive each
operation from these premises alone.

\textbf{(ii) from bounded context.} $|K| > C_n$ means the entity cannot
load all propositions simultaneously. For any problem $f$, the entity
must select which propositions to load. If selection is random or
exhaustive, inference quality $\rho(f) \to 0$ as $|K|$ grows
(Definition~\ref{def:context}). Therefore the entity must identify and
load $K^{[f]}$, the relevant subset. This is contextual retrieval.

\textbf{(i) from evidence arrival.} Propositions do not arrive in
isolation. A new proposition $p$ is grounded in existing knowledge: the
evidence for $p$ references existing propositions. If $p$ is added
without recording what it was grounded in, then when those grounds are
later revised, the entity has no mechanism to determine that $p$ is
affected. The entity must store, at insertion time, which existing
propositions $p$ depends on. This is addition with provenance.

\textbf{(iii) and (iv) from corruptibility.} Evidence can be wrong.
When proposition $p$ is invalidated by new evidence, every proposition
whose truth value depends on $p$ is potentially compromised. The entity
must determine which propositions are grounded through $p$
(dependency query) and re-evaluate or remove them transitively
(selective removal with cascade). Without these operations, a single
corrupted proposition silently invalidates an unbounded number of
dependents, and the entity has no mechanism to detect or scope the
damage.
\end{proof}

\begin{theorem}[Graph Necessity from Retention Dynamics]
\label{thm:graph-dynamics}
Let $K$ be a knowledge base supporting the maintenance operations
derived in Lemma~\ref{lem:ops}. Restated for reference:
\begin{enumerate}
\item[(i)] \textbf{Addition with provenance}: insert a proposition $p$
   together with its evidential grounds: which existing propositions
   $p$ depends on, and by what relation.
\item[(ii)] \textbf{Contextual retrieval}: given problem $f$, retrieve
   exactly $K^{[f]}$.
\item[(iii)] \textbf{Dependency query}: given proposition $p$, determine
   which propositions are grounded through $p$.
\item[(iv)] \textbf{Selective removal with cascade}: when $p$ is
   invalidated, remove $p$ and re-evaluate every proposition whose
   grounding depends on $p$, transitively, until $K$ stabilizes.
\end{enumerate}
Any representation achieving sub-linear cost on all four operations
simultaneously necessarily encodes directed adjacency between
propositions. Under bounded active context $C_n$, directed adjacency
converges to typed edges---carrying the kind of dependency on each
edge---because untyped cascade wastes context on irrelevant
re-evaluation. A directed graph with typed edges over propositions is
a semantic graph regardless of implementation.
\end{theorem}

\begin{proof}
\textbf{Part 1: Dependency queries are necessary and continuous.}
Knowledge maintenance operates continuously: every new piece of evidence
potentially shifts the status of existing propositions. When evidence
revises or invalidates a proposition $p$, the system must evaluate every
proposition whose grounding passes through $p$. This is not a batch
operation that can be deferred; acting on a proposition whose ground
has been invalidated risks incorrect inference. The dependency query
``what is grounded through $p$?'' is therefore invoked at the frequency
of evidence arrival.

\textbf{Part 2: Lower bound without adjacency.}
Consider a representation $R$ that stores $n$ propositions without
explicit directed relationships between them (a flat store, hash table,
embedding space, or any structure that does not record which propositions
depend on which). To answer the dependency query for $p$, $R$ must
examine each of the other $n - 1$ propositions to determine whether its
grounding involves $p$. Without stored adjacency, there is no mechanism
to restrict this examination: any proposition might depend on $p$, and
the only way to determine this is to check.

In any model where the representation does not store explicit
inter-proposition relationships (cell-probe, comparison-based, or any
model where dependency is not pre-computed), the cost of a single
dependency query is therefore $\Omega(n)$. Since dependency queries are
invoked at evidence-arrival frequency (Part~1), the amortized cost of
maintenance is $\Omega(n)$ per evidence event, growing linearly with
$|K|$.

\textbf{Part 3: Cascade amplifies the cost.}
Selective removal is not a single deletion. When $p$ is removed, each
proposition $q$ grounded through $p$ must be re-evaluated: does $q$ have
alternative grounds, or should it also be removed? If $q$ also fails,
its dependents must be re-evaluated in turn. This is a transitive
closure on the dependency relation.

With explicit adjacency, cascade follows edges: cost is $O(|D_p|)$ where
$|D_p|$ is the transitive dependency set. Without adjacency, each level
of cascade requires an $\Omega(n)$ scan. Total cost: $\Omega(n \cdot L)$
where $L$ is the cascade depth, versus $O(|D_p|)$ in the graph where
$|D_p| \ll n$ for typical propositions with bounded connectivity.

\textbf{Remark.} Truth maintenance systems \cite{Doyle1979,DeKleer1986}
are a direct implementation of operations~(iii) and~(iv). A TMS stores
directed justifications (which propositions ground which) and propagates
belief revision through them when any justification is invalidated,
precisely the dependency query and cascade operations this theorem
proves are necessary. The fact that AI research independently invented
TMS to solve this exact problem is empirical evidence for the theorem's
claim: without stored directed adjacency, maintenance under evidence
arrival is intractable.

\textbf{Part 4: Any sub-linear solution encodes adjacency.}
Suppose $R$ achieves $o(n)$ cost for dependency queries. Then $R$ must
store, for each proposition, information about which other propositions
depend on it; otherwise the $\Omega(n)$ lower bound of Part~2
applies. This stored information is a set of directed relationships: for
each proposition $p$, a record of the propositions that reference $p$ as
part of their grounding.

This is precisely a directed adjacency structure. The nodes are
propositions. The edges are directed dependency relations. Any
representation that stores this information, regardless of whether it
is implemented as adjacency lists, edge tables, pointers, or any other
mechanism, \emph{is} a directed graph over propositions.

\textbf{Remark on approximate methods.}
Approximate retrieval mechanisms, learned embeddings,
locality-sensitive hashing, approximate nearest-neighbor indices,
achieve sub-linear performance on operation~(ii) (contextual retrieval)
without storing explicit directed adjacency. This does not circumvent
the theorem, because the theorem requires sub-linear cost on all four
operations \emph{simultaneously}. Approximate similarity over an
embedding space can identify propositions \emph{related to} a query, but
cannot answer ``what is \emph{grounded through} $p$?'' (operation~iii)
or propagate invalidation transitively through dependency chains
(operation~iv). These operations require knowing which propositions
depend on which, not which are \emph{nearby} in a metric space.
Proximity and dependency are distinct relations: $p$ and $q$ may be
embedding-close yet inferentially independent, or embedding-distant yet
linked by a grounding chain. Any representation that achieves sub-linear
cost on dependency queries and cascade necessarily stores directed
dependency relationships, which is directed adjacency.

\textbf{Remark on the scope of ``graph.''}
The claim is not that the representation must have some structure (which
is trivially true of any non-random encoding). The claim is that the
representation must support scoped cascade under corruption
(Lemma~\ref{lem:ops}), and that the minimal structure satisfying this
requirement is directed typed adjacency. Representations that store
structure but not directed dependency, flat stores with metadata,
embedding spaces with similarity metrics, hash tables with access
counts, fail operation~(iii): they cannot answer ``what is grounded
through $p$?'' without $\Omega(n)$ scan. The graph is not defined
broadly and then shown to be necessary. The operations are derived from
premises that do not mention graphs (Lemma~\ref{lem:ops}), and the
graph is what falls out.

\textbf{Part 5: Typed edges are the convergence point under bounded
context.}
Parts~1--4 establish that directed adjacency is necessary: without it,
dependency queries cost $\Omega(n)$. This is a hard lower bound. Typed
edges are a separate claim at a different level of necessity.

A bare adjacency structure (untyped edges) achieves sub-linear
dependency queries. It is not incorrect. But when $p$ is invalidated,
untyped cascade cannot distinguish propositions \emph{inferentially
grounded} through $p$ (which require re-evaluation) from propositions
merely \emph{associated} with $p$ (which do not). Every neighbor must
be treated as potentially dependent, expanding cascade to the full
neighborhood $|\mathrm{adj}(p)|$ rather than the true dependency set
$|D_p|$.

Under bounded $C_n$, this expansion is not merely wasteful; it
consumes the same scarce resource (active context slots) that the
entity needs for inference. Each unnecessary re-evaluation during
cascade loads propositions into context that displace
problem-relevant content, degrading $\rho$ for concurrent inference.
As $|K|$ grows and $|\mathrm{adj}(p)|$ diverges from $|D_p|$, untyped
cascade consumes an increasing fraction of $C_n$ on irrelevant
re-evaluation. Typed edges, carrying the relation type (``grounds,''
``evidences,'' ``causes,'' ``supports''), restrict cascade to
dependency-bearing edges, keeping cost proportional to $|D_p|$.

The distinction matters: directed adjacency is \emph{necessary} (no
alternative achieves sub-linear maintenance); typed edges are what
directed adjacency \emph{converges to} under bounded-context
pressure. A system with untyped edges functions but pays an
increasing tax on $C_{\mathrm{free}}$ as knowledge grows. This is
the same convergence pattern---not logical necessity, but what
bounded systems converge to under $C_n$ pressure---that recurs
whenever a finite-capacity system must scale.

\textbf{Remark on empirical convergence.}
The history of knowledge representation confirms this trajectory.
Early semantic networks used untyped ``is-related-to'' links;
Woods \cite{Woods1975} demonstrated that these were semantically
ambiguous and operationally inadequate: the system could not
distinguish inheritance from association, causation from correlation.
The field converged to typed relations: RDF introduced typed
predicates \cite{BernersLee2001}, OWL distinguished object properties,
data properties, and annotation properties with transitive, symmetric,
and inverse role characteristics \cite{BaaderCalvaneseEtAl2003}.
This trajectory, starting untyped, encountering scaling failure,
converging to typed, is the progression Parts~1--5 predict.

A directed graph with typed edges over propositions is a semantic graph.
\end{proof}

\textbf{Remark (Cascade as error recovery).}
The cascade mechanism proven necessary in Parts~3--5 is not only an
efficiency requirement. It is the error recovery mechanism. When
evidence corrupts a proposition $p$, the system must identify every
proposition whose grounding passes through $p$, re-evaluate or remove
each, and scope the damage to the actual dependency set $D_p$. Without
stored typed adjacency, this scoping is impossible: the system cannot
distinguish propositions affected by the corruption from those that are
not, and must treat the entire knowledge base as potentially
contaminated. Stored typed adjacency reduces error recovery from
$\Omega(n)$ (full re-evaluation) to $O(|D_p|)$ (dependency-scoped
re-evaluation). The maintenance structure and the robustness structure
are the same structure.

\begin{corollary}[Edge Growth Dominates Node Growth]
\label{cor:edge-growth}
As $K$ matures under continuous evidence arrival, the ratio of edges to
nodes increases. Each new proposition $p_{new}$ added to $K$ arrives with
at least one edge (its evidential ground, by operation~(i) of
Theorem~\ref{thm:graph-dynamics}). As existing propositions are
cross-grounded through new evidence, new edges are added between
existing nodes without adding new nodes. Edges are the cheapest
structure in the graph: each edge is a typed pointer between existing
nodes, not new propositional content. The information density of the
graph, the ratio of relational information to propositional content,
therefore increases as $K$ matures. A mature knowledge base is
edge-dense: most of its informational value is in the relationships
between propositions, not in the propositions themselves.
\end{corollary}

% ------------------------------------------------------------------
\section{Continuous-Weight Mechanisms Are Graph Implementations}
\label{sec:attention}
% ------------------------------------------------------------------

A natural objection: modern attention mechanisms achieve contextual
retrieval without explicit graph construction. Does this constitute a
non-graph alternative?

\begin{theorem}[Contextual Retrieval Requires Adjacency Computation]
\label{thm:attention}
Any mechanism that performs contextual retrieval over $n$ propositions,
selecting which propositions are relevant to a given query, must
compute pairwise relevance between the query and stored propositions.
This computation either (a) evaluates all $n$ propositions per query at
$O(n)$ or $O(n^2)$ cost (dynamic adjacency), or (b) pre-computes and
stores relevance relationships at $O(d)$ per query where $d$ is the
local degree (stored adjacency). In case (b), the stored relationships
constitute a directed graph. Case (a) is the cost of recomputing that
graph from scratch at every query.
\end{theorem}

\begin{proof}
Contextual retrieval for problem $f$ requires identifying $K^{[f]}$
, the propositions relevant to $f$. This requires determining, for
each proposition $p \in K$, whether $p$ is relevant to $f$. Without
stored information about relevance relationships, every proposition must
be evaluated: cost $\Omega(n)$ per query.

Full self-attention is the canonical example:
$A_{ij} = \mathrm{softmax}(Q_i K_j^\top / \sqrt{d_k})$ evaluates all
$n^2$ pairwise relevances per query. This is not ``computing a graph''
in the trivial sense that any matrix can be called a graph; it is
performing the specific computation (pairwise relevance determination)
that a pre-built graph would make unnecessary. The attention matrix is
dense: all $n^2$ entries are nonzero, encoding no structural
information. The computational cost is the cost of having no stored
structure.

Pre-computed adjacency, stored information about which propositions
are relevant to which queries, reduces retrieval to $O(d)$ where $d$
is the number of stored relevant relationships. This is a directed
graph: nodes are propositions, edges are relevance relationships, and
the direction encodes which proposition informs which query context.

Sparse attention mechanisms, learned sparsity patterns, local
windows, routing functions, reduce the $O(n^2)$ cost by learning
which pairwise relevances are worth computing: they are learning which
edges to store. The KV cache materializes computed relevances for reuse
across queries. Each optimization converts dynamic adjacency computation
into stored adjacency, progressively building the graph that
Proposition~\ref{thm:graph-info} and Theorem~\ref{thm:graph-dynamics} prove is
necessary.

The result is not that ``attention is a graph'' (which would be
trivially true of any matrix) but that contextual retrieval requires
adjacency information, and the only question is whether that information
is computed dynamically ($O(n^2)$) or stored explicitly ($O(d)$). Stored
adjacency over propositions with typed relevance relationships is a
directed graph.
\end{proof}

\textbf{Remark (Scaling context windows).}
Growing $C_n$ directly does not solve the retrieval problem. For any
specific problem $f$, the relevance density $\rho(f) = |K^{[f]}| /
|\mathrm{ctx}(t)|$ (Definition~\ref{def:context}) decreases as context
fills with content not relevant to $f$, an effect documented
as ``lost in the middle'' \cite{Liu2024}. The efficient
path is not to grow the context window but to grow the encoded knowledge
accessible via stored adjacency: filling the graph, not the context
window.

% ------------------------------------------------------------------
\section{Discussion}
\label{sec:discussion}
% ------------------------------------------------------------------

The two results establish graph necessity from independent premises.
Proposition~\ref{thm:graph-info} requires only the curation regime (bounded
capacity, new propositions arriving). Theorem~\ref{thm:graph-dynamics}
requires only continuous knowledge maintenance (evidence arrives,
dependencies must be traced) and does not assume lossless preservation
; it applies equally to systems that tolerate information loss,
because the cost argument concerns maintenance operations, not
retention policy. Any system subject to both constraints,
which includes every physical system that accumulates knowledge, is
doubly committed to graph structure.

\subsection{The continuous representation objection}

The strongest objection to this paper's results is that they
assume discrete propositions, while the most capable knowledge systems
, large language models, biological neural networks, store
knowledge in continuous distributed representations that do not
decompose into identifiable propositions. If there are no discrete
propositions, there is nothing to factor
(Proposition~\ref{thm:graph-info}), nothing to query dependencies over
(Theorem~\ref{thm:graph-dynamics}), and calling an attention matrix a
``graph'' (Theorem~\ref{thm:attention}) is a mathematical relabeling,
not a structural discovery. The objection: the theorems prove graph
necessity for an idealized system that no successful architecture
actually implements.

This objection is serious and requires a substantive response. We offer
four.

\textbf{(1) Continuous representations do encode discrete
propositions.}
The continuity of the substrate does not entail that the information
content is non-propositional. Meng et al.\ \cite{Meng2022} demonstrate
that specific factual associations in GPT-class models can be localized,
edited, and ablated at identified parameter subsets: ``The Eiffel Tower
is in Paris'' has a locatable representation that can be revised to ``The Eiffel Tower is in Rome'' without retraining. Geva
et al.\ \cite{Geva2021} show that feed-forward layers in transformers
function as key-value memories, with individual neurons activating for
identifiable semantic patterns. Li et al.\ \cite{Li2024} identify discrete computational circuits
within continuous networks for mathematical reasoning.

These results establish that continuous weight matrices encode
propositional content, not approximately or metaphorically, but in a
sense precise enough to support targeted editing. The propositions are
\emph{superposed} (multiple propositions share parameters) and
\emph{distributed} (each proposition spans multiple parameters), but
they are identifiable, editable, and individually removable. Superposed
encoding is precisely the condition that makes factoring valuable: shared
structure across propositions is the input to factoring, and continuous
networks contain it abundantly.

\textbf{(2) The claim is about information content, not
implementation.}
Proposition~\ref{thm:graph-info} and Theorem~\ref{thm:graph-dynamics} prove that
the \emph{information content} of any adequate representation has graph
structure. They do not prescribe adjacency lists, edge tables, or
symbolic labels. A neural network whose weights encode directed
dependencies between propositions \emph{is} maintaining a graph in the
sense of Definition~\ref{def:knowledge}: the graph exists in the
information, not in the data structure.

The objection that ``if any structured representation counts as a
graph, the claim is trivially true'' misidentifies the contribution.
The claim is not that representations have structure (which is trivial)
but that the structure must be \emph{directed}, \emph{typed}, and
\emph{adjacency-preserving}, and that representations lacking these
specific properties lose information or pay $\Omega(n)$ maintenance
costs. Not all structure is graph structure. A flat embedding space has
structure (metric, topology) but not directed adjacency; a hash table
has structure (key-value mapping) but not typed dependencies. The
theorems identify \emph{which} structural properties are necessary, not
merely that some structure is required.

\textbf{(3) Catastrophic forgetting is the theorem's prediction for
systems without explicit adjacency.}
This is the strongest empirical evidence for the result.
Theorem~\ref{thm:graph-dynamics} proves that without stored directed
adjacency, selective removal with cascade, updating one proposition
and propagating correctly to its dependents, costs $\Omega(n)$ per
evidence event. The prediction: systems that store propositions in
continuous weights without explicit adjacency should exhibit
\emph{uncontrolled interference when knowledge is updated}.

This is precisely catastrophic forgetting
\cite{McCloskey1989,Kirkpatrick2017}: updating a neural network on new
information degrades previously learned information unpredictably.
The network cannot scope the update because it lacks the adjacency
to determine what depends on the modified parameters. Continual
learning remains an open problem in deep learning not because of
insufficient engineering but because the underlying representation
lacks the structural property the theorem identifies: directed
adjacency that enables scoped updates.

The response to the continuous representation objection is therefore
not defensive but \emph{predictive}: the theorem identifies the specific
structural deficit (missing explicit adjacency), predicts the specific
failure mode (unscoped updates causing cascading interference), and the
most successful continuous representation systems exhibit exactly that
failure mode. The objection that the theorems do not apply to neural
networks is refuted by the fact that neural networks suffer from
precisely the pathology the theorems predict for representations
without graph structure.

\textbf{(4) Knowledge graph embeddings are continuous implementations
of graph structure.}
A further data point: knowledge graph embedding methods (TransE
\cite{Bordes2013}, RotatE \cite{Sun2019}) project graph structure into
continuous vector spaces. The embeddings are continuous, but the
information content they encode, which entities are related to which,
by what relation, is directed typed adjacency. Inference tasks such
as link prediction require reconstructing adjacency from the embedding
space. These methods do not eliminate graph structure; they implement it
in a continuous substrate, which is exactly the relationship
Proposition~\ref{thm:graph-info} and Theorem~\ref{thm:graph-dynamics} predict between
information content (graph) and implementation (free parameter).

\subsection{Scope and non-claims}

The theorems prove that the \emph{information content} of the
representation must be a directed graph. They do not prescribe a
specific graph formalism (property graph, RDF, Bayesian network) or
implementation (adjacency lists, sparse matrices, attention masks). The
graph is the abstract structure. The substrate is a free parameter.

\textbf{Implications for AI architecture.}
Theorem~\ref{thm:attention} reframes the relationship between
graph-based and neural approaches: they are not competing paradigms but
different implementations of the same structural necessity. Systems that
compute adjacency dynamically (full attention) pay $O(n^2)$; systems
that store it explicitly (graphs, sparse attention, KV caches) pay
$O(d)$. The trend in architecture design, from dense attention toward
sparse, structured, and cached retrieval, is convergence toward
explicit graph structure, driven by the same cost pressures the theorems
identify.

The result suggests that hybrid architectures combining explicit
knowledge graphs with neural retrieval are not an engineering
convenience but a response to a structural constraint: as $|K|$ grows,
systems must encode adjacency to maintain inference quality under bounded
context.

\textbf{Free context and the purpose of graph structure.}
The graph structure derived in this paper exists to serve a single
resource: \emph{free context}, the portion of $C_n$ not consumed by
maintenance of existing knowledge and available for novel problems.
Formally, $C_{\mathrm{free}}(t) = C_n - |\mathrm{ctx}(t)|$,
where $|\mathrm{ctx}(t)|$ is the number of context slots occupied by
currently active propositions. Factoring
(Proposition~\ref{thm:graph-info}) increases $C_{\mathrm{free}}$ by
compressing shared structure into fewer propositions; stored adjacency
(Theorem~\ref{thm:graph-dynamics}) increases $C_{\mathrm{free}}$ by
enabling targeted retrieval, loading only the $d$ relevant propositions
rather than scanning all $n$, so that fewer context slots are consumed
by irrelevant material during inference. Every result in this paper is, at root,
a mechanism for maximizing $C_{\mathrm{free}}$: the context capacity
available for novel problems after existing-knowledge maintenance.

\textbf{Implications for scaling.}
The scaling remark following Theorem~\ref{thm:attention} challenges the
prevailing approach of growing context windows as the primary path to
more capable inference.
Under bounded context, larger windows with undifferentiated content
\emph{reduce} inference quality for any specific problem. The
alternative, encoding knowledge into a graph and retrieving the
relevant subgraph per problem, maintains $\rho = 1$ regardless of
total knowledge size. This is the difference between scaling the window
and scaling the knowledge.

\subsection{Falsifiable predictions}

The theorems yield seven predictions that are empirically testable. Each
states a consequence of the formal results and the observation that
would falsify it.

\textbf{1. Catastrophic forgetting is structural, not parametric.}
Theorem~\ref{thm:graph-dynamics} predicts that systems storing
propositions in continuous weights without explicit adjacency exhibit
uncontrolled interference when knowledge is updated. This is
catastrophic forgetting. The prediction: continual learning will not be
solved by regularization, replay, or parameter isolation alone. Any
solution that works will encode directed adjacency over propositions,
whether explicitly or implicitly. A system that solves continual
learning without adjacency would falsify Theorem~\ref{thm:graph-dynamics}.

\textbf{2. Architectural convergence toward stored adjacency.}
Theorem~\ref{thm:attention} proves that contextual retrieval without
stored adjacency costs $O(n^2)$; with stored adjacency, $O(d)$. The
prediction: the dominant trajectory in neural architecture design will
move from dense attention toward sparse, structured, cached retrieval,
progressively building explicit adjacency. Sparse attention patterns,
learned routing, KV caches, and retrieval-augmented generation are early
instances. More specifically, learned sparsity patterns will
increasingly resemble stored dependency structure: current routing in
mixture-of-experts is content-based; the prediction is that effective
routing will become topology-aware, because that is what the cost
pressure selects for. A system that achieves sub-linear contextual
retrieval without encoding any form of directed adjacency would falsify
Theorem~\ref{thm:attention}.

\textbf{3. Context-optimizing architectures will outcompete
context-expanding architectures on cost.}
Current large language models maximize functional context by expanding
the context window. Theorem~\ref{thm:attention} shows this cannot
scale: relevance density $\rho$ degrades as undifferentiated content
fills the window, and the retrieval cost remains $O(n^2)$. An
architecture that accrues structured knowledge and retrieves the
relevant subgraph per problem maintains $\rho = 1$ at $O(d)$ cost
regardless of total knowledge size. The prediction: as task complexity
and knowledge volume grow, systems that optimize context (structured
accrual + selective retrieval) will undercut systems that expand context
(larger windows + denser packing) on cost per correct answer, and
competitive pressure will drive convergence toward the former. A
context-expanding architecture that maintains cost parity with
structured retrieval at arbitrary knowledge scale would falsify
Theorem~\ref{thm:attention}.

\textbf{4. Context window scaling yields diminishing returns.}
For any problem $f$, relevance density
$\rho(f) = |K^{[f]}| / |\mathrm{ctx}(t)|$ decreases as context fills
with content irrelevant to $f$ (Definition~\ref{def:context}). The
prediction: increasing context window size while holding knowledge
structure constant will show diminishing returns in task performance,
and eventually negative returns for retrieval-sensitive tasks. The
efficient scaling path is growing encoded knowledge accessible via
stored adjacency, not growing the context window.

\textbf{5. Retrieval-augmented systems outperform context stuffing.}
Direct consequence of Theorem~\ref{thm:attention}. Retrieving the
relevant subgraph $K^{[f]}$ and loading it into context maintains
$\rho = 1$ regardless of total knowledge size. Stuffing all available
context indiscriminately degrades $\rho$. The prediction: for any fixed
context window, retrieval-augmented architectures will outperform
context-stuffing architectures on knowledge-intensive tasks, with the
gap widening as total knowledge grows.

\textbf{6. Mature knowledge bases are edge-dense.}
Corollary~\ref{cor:edge-growth} predicts that the ratio of edges to
nodes increases as $K$ matures under continuous evidence arrival. The
prediction: empirical knowledge graphs (Wikidata, biomedical
ontologies, institutional knowledge bases) will show increasing
edge-to-node ratios over time, and the informational value of the graph
(measured by downstream task performance) will correlate more strongly
with edge density than with node count.

\textbf{7. Error recovery cost scales with adjacency quality.}
The cascade-as-error-recovery remark predicts that the cost of
recovering from corrupted evidence scales as $O(|D_p|)$ with stored
typed adjacency and $\Omega(n)$ without it. The prediction: systems
with explicit dependency tracking will recover from misinformation or
evidence retraction faster and with less collateral damage than systems
without it. Knowledge graphs that support typed cascade will outperform
flat stores on correction tasks.

\textbf{Biological evidence.}
Biological neural networks maintain sparse, structured connectivity
rather than dense all-to-all connections. Working memory is bounded
($C_n \approx 7 \pm 2$) while long-term knowledge grows without
apparent limit \cite{Miller1956}. Retrieval is associative and
context-driven, consistent with graph traversal from a query node
rather than exhaustive scan. These observations are consistent with
the theorems' predictions, though the biological case requires
caution: the theorems assume propositional decomposition, and the
degree to which biological representations admit such decomposition
remains an open empirical question.

% ------------------------------------------------------------------
\section{Conclusion}
\label{sec:conclusion}
% ------------------------------------------------------------------

We have shown that graph structure over propositions is not a design
choice but a necessary consequence of two independently sufficient
constraints: information preservation under bounded capacity, and
sub-linear knowledge maintenance under continuous evidence arrival.
Continuous-weight mechanisms such as transformer self-attention implement
graph structure dynamically rather than constituting alternatives to it.
The result applies to any system that accumulates
propositions under bounded active context and must preserve information:
such a system necessarily maintains a directed graph with typed edges, whether that
graph is realized as explicit edges, learned sparse attention, synaptic
connectivity, or any other medium.

\subsection*{Computational consequence}

A proven structural result eliminates design space permanently.
If graph structure is necessary, then any system that discovers
graph-like organization from data is rediscovering what the proof
already establishes---paying the discovery cost on every training
run for structure that was already known. The proof pays that cost
once.

The practical consequence is direct. A learning system operating
over $n$ unstructured features must evaluate $O(n^2)$ pairwise
interactions to discover which ones carry signal. A system informed
by the proven structure evaluates only the $O(k^2)$ interactions
the theory identifies as meaningful, where $k \ll n$. The
compression ratio $(n/k)^2$ is not an approximation; it is the
computational dividend of a proof. The proven structure does not
need to be rediscovered. It does not degrade. It does not require
retraining. Each theorem in this series eliminates a region of
the search space that no system---biological or artificial---needs
to explore again.

\section*{Acknowledgments}

The author used Claude (Anthropic) for drafting assistance, literature review, and \LaTeX{} preparation. All theoretical content, proofs, and arguments are the author's own.

\bibliographystyle{plain}
\begin{thebibliography}{99}

\bibitem{Arp2015}
R.~Arp, B.~Smith, and A.~D. Spear.
\newblock \emph{Building Ontologies with Basic Formal Ontology}.
\newblock MIT Press, 2015.

\bibitem{BaaderCalvaneseEtAl2003}
F.~Baader, D.~Calvanese, D.~L. McGuinness, D.~Nardi, and P.~F.
  Patel-Schneider, editors.
\newblock \emph{The Description Logic Handbook: Theory, Implementation, and
  Applications}.
\newblock Cambridge University Press, 2003.

\bibitem{BernersLee2001}
T.~Berners-Lee, J.~Hendler, and O.~Lassila.
\newblock The semantic web.
\newblock \emph{Scientific American}, 284(5):34--43, 2001.

\bibitem{Bordes2013}
A.~Bordes, N.~Usunier, A.~Garcia-Dur{\'a}n, J.~Weston, and
  O.~Yakhnenko.
\newblock Translating embeddings for modeling multi-relational data.
\newblock In \emph{Advances in Neural Information Processing Systems},
  volume~26, pages 2787--2795, 2013.

\bibitem{Codd1970}
E.~F. Codd.
\newblock A relational model of data for large shared data banks.
\newblock \emph{Communications of the ACM}, 13(6):377--387, 1970.

\bibitem{Cowan2001}
N.~Cowan.
\newblock The magical number 4 in short-term memory: A reconsideration of
  mental storage capacity.
\newblock \emph{Behavioral and Brain Sciences}, 24(1):87--114, 2001.

\bibitem{Cover2006}
T.~M. Cover and J.~A. Thomas.
\newblock \emph{Elements of Information Theory}.
\newblock Wiley-Interscience, 2nd edition, 2006.

\bibitem{DeKleer1986}
J.~de~Kleer.
\newblock An assumption-based {TMS}.
\newblock \emph{Artificial Intelligence}, 28(2):127--162, 1986.

\bibitem{Doyle1979}
J.~Doyle.
\newblock A truth maintenance system.
\newblock \emph{Artificial Intelligence}, 12(3):231--272, 1979.


\bibitem{Geva2021}
M.~Geva, R.~Schuster, J.~Berant, and O.~Levy.
\newblock Transformer feed-forward layers are key-value memories.
\newblock In \emph{Proceedings of the 2021 Conference on Empirical Methods in
  Natural Language Processing}, pages 5484--5495, 2021.

\bibitem{Guarino1998}
N.~Guarino.
\newblock Formal ontology and information systems.
\newblock In \emph{Proceedings of FOIS'98}, pages 3--15. IOS Press, 1998.

\bibitem{Hogan2021}
A.~Hogan, E.~Blomqvist, M.~Cochez, C.~d'Amato, G.~de~Melo, C.~Gutierrez,
  S.~Kirrane, J.~E.~L. Gayo, R.~Navigli, S.~Neumaier, A.-C. Ngonga~Ngomo,
  A.~Polleres, S.~M. Rashid, A.~Rula, L.~Schmelzeisen, J.~Sequeda, S.~Staab,
  and A.~Zimmermann.
\newblock Knowledge graphs.
\newblock \emph{ACM Computing Surveys}, 54(4):1--37, 2021.

\bibitem{Kirkpatrick2017}
J.~Kirkpatrick, R.~Pascanu, N.~Rabinowitz, J.~Veness, G.~Desjardins,
  A.~A. Rusu, K.~Milan, J.~Quan, T.~Ramalho, A.~Grabska-Barwinska,
  D.~Hassabis, C.~Clopath, D.~Kumaran, and R.~Hadsell.
\newblock Overcoming catastrophic forgetting in neural networks.
\newblock \emph{Proceedings of the National Academy of Sciences},
  114(13):3521--3526, 2017.

\bibitem{Li2024}
K.~Li, O.~Press, and S.~Arora.
\newblock How does {GPT}-2 compute greater-than?: Interpreting mathematical
  abilities in a pre-trained language model.
\newblock In \emph{Advances in Neural Information Processing Systems}, 2024.

\bibitem{Masolo2003}
C.~Masolo, S.~Borgo, A.~Gangemi, N.~Guarino, and A.~Oltramari.
\newblock {WonderWeb} deliverable {D18}: Ontology library (final).
\newblock Technical report, ISTC-CNR, 2003.

\bibitem{Liu2024}
N.~F. Liu, K.~Lin, J.~Hewitt, A.~Paranjape, M.~Bevilacqua, F.~Petroni, and
  P.~Liang.
\newblock Lost in the middle: How language models use long contexts.
\newblock \emph{Transactions of the Association for Computational Linguistics},
  12:157--173, 2024.

\bibitem{McCloskey1989}
M.~McCloskey and N.~J. Cohen.
\newblock Catastrophic interference in connectionist networks: The sequential
  learning problem.
\newblock In G.~H. Bower, editor, \emph{Psychology of Learning and Motivation},
  volume~24, pages 109--165. Academic Press, 1989.

\bibitem{Meng2022}
K.~Meng, D.~Bau, A.~Andonian, and Y.~Belinkov.
\newblock Locating and editing factual associations in {GPT}.
\newblock In \emph{Advances in Neural Information Processing Systems},
  volume~35, pages 17359--17372, 2022.

\bibitem{Miller1956}
G.~A. Miller.
\newblock The magical number seven, plus or minus two: Some limits on our
  capacity for processing information.
\newblock \emph{Psychological Review}, 63(2):81--97, 1956.

\bibitem{Pearl1988}
J.~Pearl.
\newblock \emph{Probabilistic Reasoning in Intelligent Systems: Networks of
  Plausible Inference}.
\newblock Morgan Kaufmann, 1988.

\bibitem{Shannon1948}
C.~E. Shannon.
\newblock A mathematical theory of communication.
\newblock \emph{Bell System Technical Journal}, 27(3):379--423, 1948.

\bibitem{Simon1955}
H.~A. Simon.
\newblock A behavioral model of rational choice.
\newblock \emph{The Quarterly Journal of Economics}, 69(1):99--118, 1955.

\bibitem{Sun2019}
Z.~Sun, Z.-H. Deng, J.-Y. Nie, and J.~Tang.
\newblock {RotatE}: Knowledge graph embedding by relational rotation in complex
  space.
\newblock In \emph{Proceedings of the International Conference on Learning
  Representations}, 2019.

\bibitem{Tishby2000}
N.~Tishby, F.~C. Pereira, and W.~Bialek.
\newblock The information bottleneck method.
\newblock In \emph{Proceedings of the 37th Annual Allerton Conference on
  Communication, Control, and Computing}, pages 368--377, 2000.

\bibitem{Woods1975}
W.~A. Woods.
\newblock What's in a link: Foundations for semantic networks.
\newblock In D.~G. Bobrow and A.~Collins, editors, \emph{Representation and
  Understanding: Studies in Cognitive Science}, pages 35--82. Academic Press,
  1975.

\bibitem{McCarthy2026e}
P.~D. McCarthy.
\newblock Coupling channels, severance hierarchy, and the prognostic
structure of cancer: an empirical test of bounded-context theory.
\newblock \emph{SSRN preprint}, 2026.

\end{thebibliography}

\end{document}
